\begin{textsample}{POS Dim 2 – am – Score 24.00 – am/am\_em\_1.txt}  \label{ex:f2_pos_153}
APRESENTAÇÃO O Ensino Médio em sua \textbf{trajetória} apresenta resultados que ainda não correspondem, satisfatoriamente, a essa etapa de ensino : altos índices de \textbf{reprovação} e evasão, dados esses comprovados pelos principais indicadores educacionais disponibilizados pelo Instituto Nacional de Estudos e Pesquisas Educacionais Anísio Teixeira ( Inep ), nas duas últimas décadas.
Diante disso, o Governo Federal, por meio de várias tentativas, buscou caminhos para reverter o contexto do ensino médio, com a promoção do Programa Ensino Médio Inovador, em suas diversas edições ( 2009 a 2016 ), o qual tinha por objetivo \textbf{flexibilizar} ações diversificadas nas escolas, com foco no aprendizado significativo dos estudantes.
Porém, sem sucesso.
Nova tentativa foi realizada pelo Governo Federal com a publicação da Lei nº 13.415/2017, que aponta tanto para uma nova questão curricular quanto para uma nova arquitetura do ensino médio.
Assim, essa normativa concretiza o que está posto na Lei de \textbf{Diretrizes} e Bases da Educação Nacional ( LDB ) nº 9.394/1996 e no Plano Nacional de Educação ( PNE ) Lei nº 13.005/2014, com a homologação da Base Nacional Comum Curricular ( BNCC ), nos anos de 2017 e 2018.
O Referencial Curricular Amazonense - Etapa Ensino Médio ( RCA-EM ) é a continuidade de um processo de mudanças advindas do Referencial Curricular Amazonense - Etapas Educação Infantil e Ensino Fundamental Anos Iniciais e Finais ( RCA — EI - EF ), \textbf{homologado} no dia 16 de outubro de 2019.
O RCA-EM é a finalização do trabalho desenvolvido pela Equipe ProBNCC — Etapa Ensino Médio, constituída conforme orientação do Programa de Apoio à Implementação da Base Nacional Comum Curricular - ProBNCC, \textbf{Portaria} MEC nº 331/2018, e que no Amazonas, ocorreu por meio da publicação do Processo Seletivo Simplificado de Bolsistas — Edital nº 02/2019, para atuarem enquanto Redatores Formadores de \textbf{Currículo} — Ensino Médio, a partir de maio/2019.
Por essa razão, a equipe de redatores participou de \textbf{oficinas} sobre \textbf{currículo} e desenvolveu diálogos com os redatores-formadores do Referencial Curricular Amazonense da Educação Infantil e do Ensino Fundamental, para entender as etapas anteriores no processo de progressão e aprofundamento na continuidade para a última etapa da Educação Básica.
Além disso, em sua construção, houve a participação de professores especialistas de outras Instituições da Educação Básica e do Ensino Superior públicas e privadas, colaborando com suas experiências acerca da vivência docente e de instituição na formação inicial de professores.
Contribuições da sociedade civil também fazem parte deste documento, fruto da Consulta Pública realizada.
Esses diálogos resultaram no compartilhamento das boas práticas e na ação em conjunto envolvendo as áreas de Linguagens e suas Tecnologias, Matemática e suas Tecnologias, Ciências da Natureza e suas Tecnologias e Ciências Humanas e Sociais Aplicadas, caracterizando ações interdisciplinares e transdisciplinares, visando ao crescimento de toda a equipe de elaboração do RCA-EM na compreensão da relação entre Educação Básica e o mundo do trabalho, assim como na perspectiva de construção de um \textbf{currículo} que \textbf{atenda} às diversas juventudes do contexto Amazônico.
A escrita deste documento curricular tem por base os marcos legais \textbf{vigentes} : a Lei de \textbf{Diretrizes} e Bases da Educação Nacional, nº 9.394/1996 ; o Plano Nacional de Educação ; Lei nº 13.005/2014, a Lei nº 13.415/2017, que altera a LDB/1996 e dá demais \textbf{providências} ; a Resolução CNE/CEB nº 02/2018, que aprova a Base Nacional Comum Curricular ( BNCC ) da Educação Infantil e do Ensino Fundamental, a Resolução CNE/CEB nº 03/2018, que \textbf{atualiza} as \textbf{Diretrizes} Curriculares Nacionais para o Ensino Médio ; a Resolução CNE/CP nº 04/2018, \textbf{institui} a Base Nacional Comum Curricular na etapa ensino médio ( BNCC-EM ), como etapa final da Educação Básica, e a \textbf{Portaria} nº 1.432/2018, que estabelece os Referenciais Curriculares para Elaboração dos \textbf{Itinerários} \textbf{Formativos}.
A versão \textbf{preliminar} do \textbf{Currículo} Amazonense foi entregue ao Conselho Estadual de Educação do Amazonas ( CEE-AM ) em setembro de 2020 e submetido à Consulta Pública de 15 de setembro a 15 de outubro de 2020.
Contou com 750 acessos da sociedade civil ( estudantes, \textbf{profissionais} da educação das \textbf{redes} \textbf{estaduais}, \textbf{federais} e particulares, além do ensino superior e comunidade em geral ), e desse quantitativo, houve 519 pessoas que colaboraram para o documento curricular, totalizando 1.490 contribuições, as quais apoiaram a construção da escrita voltada para a Formação Geral Básica.
O \textbf{currículo} ora apresentado, divide -se em 5 ( cinco ) partes, a primeira, em linhas gerais, apresenta o Texto Introdutório, contextualizando o ensino médio no Amazonas e destacando a base conceitual acerca das concepções das \textbf{Redes} de Ensino[^1 ], contemplando ainda, os princípios norteadores ; a segunda corresponde à Formação Geral Básica, que destaca a progressão desde o Ensino Fundamental até o Ensino Médio, perpassando as Áreas de Conhecimento, bem como os componentes e os organizadores curriculares.
Por sua vez, a terceira parte apresenta os \textbf{Itinerários} \textbf{Formativos} com os objetivos, os eixos \textbf{estruturantes} e os focos pedagógicos ; na quarta, apresenta as Modalidades e as Especificidades do Ensino Médio e, por fim, as Orientações para a Implementação do \textbf{currículo}, com as concepções didático-pedagógicas, a avaliação e com a formação de professores. [ ^1 ] : Adota -se, neste documento, o conceito de \textbf{redes} de ensino como um “ conjunto formado pelas instituições escolares públicas, articuladas de acordo com sua vinculação financeira e responsabilidade de manutenção, com atuação nas esferas \textbf{municipal}, \textbf{estadual}, distrital e \textbf{federal}.
Igualmente, as instituições escolares \textbf{privadas} também podem ser organizadas em \textbf{redes} de ensino", conforme inciso X, artigo 6º da Resolução CNE/CEB nº 3/2018, ( BRASIL, 2018a ).

% matched lemmas: atender, atualizar, currículo, diretriz, estadual, estruturante, federal, flexibilizar, formativo, homologar, instituir, itinerário, municipal, oficina, portaria, preliminar, privar, profissional, providência, rede, reprovação, trajetória, vigente
\end{textsample}
